\section{Conclusion}
\label{sec:conclusion}

This review asked how well deep-learning-based waste detection generalizes beyond its training environment, synthesizing 50 studies across platforms (Section~\ref{sec:datasources}), approaches (Section~\ref{sec:approaches}), architectures (Section~\ref{sec:architectures}), datasets (Section~\ref{sec:datasets}), training strategies (Section~\ref{sec:training}), and evaluation practice (Section~\ref{sec:evaluation}). Three findings stand out. First, transfer works within bounds: proximate rivers, similar sites, and matched spectra transfer reliably, while sediment shifts, land-to-water jumps, and marine-to-river leaps do not. Second, the field is solving annotation scarcity faster than it is solving measurement completeness: open datasets, self-supervision, synthetic supplementation, and adaptation recover data efficiency, yet submerged, transparent, and entangled litter still defeats quantification. Third, evaluation lags capability: detection scores systematically overstate counting and flux reliability.

The review itself is limited by seed-anchored (not exhaustive) collection, single-agent screening, English-only sources, a 2019--2025 window, and provisional rows quarantined from strong claims. On the base study \cite{Pati2025}, this review supports its cross-river transfer finding across many independent replications while qualifying its training-strategy finding as architecture- and task-dependent. Its contribution is the first generalization-centered synthesis of the field with explicit evidence grading, giving practitioners a transferability map and researchers a gap inventory with actionable next steps.
